Both experiments use the frozen code (\code{autose-v1}, one crash fix in \code{v1.1}), the same five issues, the same hidden tests and the same gates (\code{experiment/variants.md}, \code{experiment/protocol.md}). Reruns after infrastructure faults are listed in \code{results/runs.md}; none of the results below comes from a run that stopped because of quota or an outage.

\textbf{E1, splitting the work (\adr{001}).} The structure divides stages 1-4 into roles of a planner, an implementer and a tester, each role having their specific inputs. The modification gives the same agent (the same model, \code{gpt-oss-120b}) all of the tasks in a single conversation with the same gates applied to its final product and a maximum of three attempts. The results of round 1 were compared (Table~\ref{tab:e1}): the pipeline has opened pull requests on issue \#4, \#6, and \#7 (hidden tests 4/4, 4/4 and 2/4 respectively, no regressions) and was stopped at stages 4 (\#5) and 2 (\#8), whereas the agent produced \emph{no} pull request on all of the five issues. Out of the 15 attempts the agent was able to pass spec, design, and implementation tests five times, but failed the test gate every time (no test named \code{test\_ac<N>} was found). In eight attempts it lost the required form of the answer (malformed or repeated tool calls, a request over the context limit). The use of tokens remains the same (from 102,000 to 218,000 for the single agent and from 120,000 to 228,000 for pipeline), so the difference is not just a matter of budget. Our interpretation is that one conversation within a 7000-token context cannot combine requirements, design, code and testing, and there is nothing that provides hints about which component failed; the pipeline gives each stage a limited context and explicit forward and backward edges. This does not show that splitting is better than a single agent with strong models, and the single agent has no stage-level feedback, so the experiment tests splitting together with gating. The result supports \adr{001}; we did not change the design because of it.

\textbf{E2, reviewer independence (\adr{008}).} The final reviewer was run three times on each of the five seeded pull requests, for each variant (Table~\ref{tab:e2}). The independent reviewer (Nemotron, own job, issue plus diff plus repository) approved all six reviews of the correct pull requests and caught the defect in eight of its nine reviews of the wrong ones; the ninth gave no valid verdict and counts as no review. It found the bugs by running snippets of its own (on pull request \#16 it showed that \code{startswith('al')} returned the results of an earlier \code{startswith('Al')}). The same-model reviewer (the builder's \code{gpt-oss-120b}, fresh context) approved eight of nine wrong pull requests: it ran the test suite once, read the pull request's own tests and approved, using about 3,400 tokens against 54,000. Our evidence rule did not help in this case, because it verifies that evidence exists, not that it could have found a defect. The result supports \adr{008} (the choice of a different model and of a verification habit matter, not the separate job alone).

\textbf{E2b, a larger reviewer (bonus).} Since the independent reviewer's first model has 12 billion active parameters while the builder's model has only about 5 billion, we asked whether the results depend on the size of the reviewer model. We conducted another test with the Nemotron 550 billion (55 billion active parameters) as the first tier. It produced 15 valid reviews, caught all nine wrong pull requests and approved all six correct ones, at about twice the tokens (105,000 per review). Both the reviewers performed considerably better than the same-model version; however, the model with larger parameters performed better in terms of reliability (no invalid review). As a result, we had to modify the design of the method \emph{after} these runs were conducted (the reviewer chain now starts with the larger model; see \adr{008}, \code{autose-v1.2}). The recorded pipeline reviews were made by the smaller first tier, as frozen.
